\subsection{Simultaneous absolute repairs}

Write $\mathcal S_k=s\operatorname{Mod}_k$ and write
$\underline{k}$ for its constant object. Cochain degree
below is the form degree; normalization has homological
degree. We first record the two comparison facts needed
to distinguish an actual obstruction from its possible
group.

For every $Y\in\mathcal S_k$ there is a natural equivalence
\begin{equation}
 R\operatorname{Hom}_{\mathcal S_k}(Y,\underline{k})
 \simeq R\operatorname{Hom}_k(NY,k).
 \tag{B-C1}
\end{equation}
Here a normalized term in homological degree $n$ is
placed in cochain degree $-n$ when forming the right
side. In particular all the groups
$\operatorname{Ext}^t_{\mathcal S_k}(Y,\underline{k})$
vanish if $NY$ is acyclic, without a flatness
assumption on $Y$.

Indeed $H_0:\operatorname{Ch}_{\geq0}(k)\to
\operatorname{Mod}_k$ is left adjoint to the exact
constant-complex functor. Projectives in the complex
category are retracts of sums of the identity disks
$D_n(P)$, $n\geq1$, and $P[0]$, with $P$ projective.
Their $H_0$ is projective and their higher homology
vanishes. If $P_\bullet\to NY$ is a projective
resolution in the abelian category of complexes, the
first-quadrant double complex $P_\bullet$ has total
complex quasi-isomorphic both to $NY$ and to
$H_0(P_\bullet)$. Thus $LH_0(Y)\simeq NY$.
Applying the derived adjunction proves (B-C1).

We also need a bounded comparison. Let $C$ be a
cochain complex in $\mathcal S_k$, supported in
degrees $0,\ldots,b$, whose normalized terms
$NC^i$ are degreewise projective. Let $A$ be
nonnegative in cochain degree and suppose that
$NA^j$ is acyclic for $0\leq j\leq b$.
Then
\begin{equation}
 H^0\operatorname{Tot}\operatorname{Hom}_{\mathcal S_k}(C,A)
 \ \cong\
 H^0R\operatorname{Hom}_{\mathcal S_k}(C,A).
 \tag{B-C2}
\end{equation}
To prove this, first observe that
$\operatorname{Ext}^{t>0}_{\mathcal S_k}(C^i,A^j)=0$
in that range. An extension of the corresponding
normalized complexes splits degreewise. Its class
is a chain-homotopy class into a shift of the
acyclic target, and is zero because the source is
bounded below and degreewise projective. A
projective-object resolution of $C^i$ has
degreewise-projective syzygies: every termwise
epimorphism onto a projective module splits.
Dimension shifting proves the assertion for all $t>0$.

In the hyper-Ext comparison the entries have total
degree $t+j-i$, where $0\leq i\leq b$.
An entry with $t>0$ and total degree at most one
has $j\leq b$. All these entries therefore vanish.
The ordinary Hom complex is the $t=0$ edge, giving
(B-C2). This is a bounded comparison: $A^{b+1}$
can occur in the ordinary differential, but no
horizontal Ext-acyclicity assumption on that row
is required. The filtrations converge because
$i$ is bounded and $j,t$ are nonnegative.

If in addition $\underline{k}\to A$ is an
isomorphism on vertical $H^0$ and $H^j(A)=0$
for $1\leq j\leq b$, then
\begin{equation}
 H^0R\operatorname{Hom}_{\mathcal S_k}(C,\underline{k})
 \ \cong\
 H^0R\operatorname{Hom}_{\mathcal S_k}(C,A).
 \tag{B-C3}
\end{equation}
The cone of $\underline{k}\to A$ has vertical
cohomology only in degrees at least $b+1$.
The source is supported in degrees at most $b$,
so its derived Hom into that cone has no
cohomology in degrees at most zero. This
proves the assertion. This last argument
does not require projectivity of the terms
of $C$.

\begin{proposition}[C22: \status{Proved}; the actual second quadratic repair]
\label{sa08:C22}
Let $P=\Gamma^2K(\Z,2)$ and let
$z:P\to A^{2w,3}$ be any prescribed simplicial
linear map. Suppose that in weight $4w$ the
target is a vertical resolution of its
coefficient $\underline{\Z}$ through degree
eight, and that its rows zero through eight
are horizontally acyclic. Then the quadratic
part of the odd-to-even disk $z\mapsto dz$
admits a strict simultaneous simplicial repair.
The prescribed map $z$ is retained. Thus the
nonzero possible group in Statement~\ref{sa08:C20}
does not force a nonzero actual obstruction.
\end{proposition}
\begin{proof}
The quadratic disk is
\[
 C=\bigl(\Lambda^2P\longrightarrow P\otimes P
             \longrightarrow\Gamma^2P\bigr)
\]
in form degrees $6,7,8$. For a free abelian
module, the multiplication map onto the
divided square has cokernel $P/2$, the
antisymmetric tensors give its kernel,
and the exterior map is injective.
The identification of the cokernel is natural:
$\gamma_2(a)\mapsto a\bmod2$ and mixed
products map to zero. Consequently
$C\simeq(P/2)[-8]$ as a cochain complex
in $\mathcal S_{\Z}$.

The complex $N(P/2)$ is supported in normalized
degrees two through four and has terms of
projective dimension at most one over $\Z$.
Its projective dimension in the abelian
category of nonnegative complexes is at
most five. For completeness, a sphere
complex $V[n]$ has projective dimension
at most $n+\operatorname{pd}_{\Z}V$:
use the identity-disk resolutions and
then a projective resolution of $V$.
Filter a bounded complex by its terms
to get the asserted bound. Hence
\[
 \operatorname{Ext}^8_{\mathcal S_{\Z}}
       (P/2,\underline{\Z})=0.
\]
It also follows directly from (B-C1) and
the same normalized degree bound.

The given algebra map determines a cochain
map $f:C\to A^{4w,*}$. All normalized terms
of $C$ are projective. Applying (B-C2)
and (B-C3), with $b=8$, gives
\[
 H^0\operatorname{Tot}\operatorname{Hom}_{\mathcal S_{\Z}}
       (C,A^{4w,*})
 \cong
 \operatorname{Ext}^8_{\mathcal S_{\Z}}
       (P/2,\underline{\Z})=0.
\]
Thus there is an actual simplicial linear
cochain homotopy $h$, with $dh+hd=f$.
Adjoin generators $sc$ in form degree
$|c|-1$ and parameter arity two, for
$c\in C$, with
\[
 d(sc)=c-s(dc).
\]
Their map to the target is $sc\mapsto h(c)$.
The formulas give $d^2=0$, and the free
PD algebra property extends the map.
The corrected quadratic complex is the
identity cone on $C$ and is contractible.
All primitive choices have been made
together; none changes the original $z$.
\end{proof}

\begin{proposition}[C23: \status{Proved}; absolute fixed-arity disk repair]
\label{sa08:C23}
Let $q\geq1$ be odd, let $m\geq2$, and
let $P$ be a simplicial $k$-module whose
normalized terms are projective. Suppose
$x:P\to A^{w,q}$ is simplicial linear.
Set $b=m(q+1)$, and assume:
\begin{enumerate}
\item the row $A^{w,q}$ is degreewise
flat and horizontally acyclic;
\item in weight $mw$, the vertical
augmentation $\underline{k}\to A^{mw,*}$
is an isomorphism on $H^0$, and
$H^j(A^{mw,*})=0$ for $1\leq j\leq b$;
\item the rows $A^{mw,j}$ are
horizontally acyclic for $0\leq j\leq b$.
\end{enumerate}
Then the map from the free odd-to-even
disk $x\mapsto dx$ admits a strict
simultaneous simplicial repair of its
parameter-arity-$m$ complex. One may
freely adjoin a shifted copy of that
complex with identity-cone differential.
All added generators have positive
form degree.
\end{proposition}
\begin{proof}
Write $D_m(P)$ for the specified arity
of the free disk. Its terms are
\[
 \Lambda^aP\otimes\Gamma^cP,\qquad a+c=m,
 \quad\text{in form degree }mq+c.
\]
They have degreewise-projective normalized
terms: exterior and divided powers preserve
projectives by their free-module bases and
preservation of retracts, and normalization
is a direct summand of the simplicial term.
The complex is bounded above by $b$.

Put $U=A^{w,q}$. The actual disk map
factors as a map of cochain complexes
\[
 D_m(P)\longrightarrow D_m(U)
       \longrightarrow A^{mw,*}.
\]
In the second arrow the formal odd
generator is sent to its underlying
element of $U$, and its even
differential is sent to the actual
differential in $A$.
Every horizontal term of $D_m(U)$
is acyclic. Indeed a positive divided
or exterior power of $U$ is flat
and tensor-universally acyclic by
Statement~\ref{sa08:C05}. The displayed
mixed tensor terms are therefore
acyclic as well, including the cases
where one factor has exponent zero.

By (B-C1), the derived Hom of each
of these terms into $\underline{k}$
vanishes. Their finite cochain
filtration gives
$R\operatorname{Hom}_{\mathcal S_k}
 (D_m(U),\underline{k})\simeq0$.
Applying (B-C3) to $D_m(U)$ proves
that its map to $A^{mw,*}$ is zero
in the derived category. Hence the
actual map $D_m(P)\to A^{mw,*}$
is zero there as well. By (B-C2),
it has an actual simplicial linear
cochain homotopy to zero.

Adjoin $sD_m(P)$ with
$d(sc)=c-s(dc)$ and use the homotopy
as its image in $A$. This extends
the prescribed original disk map.
The smallest new form degree is
$mq-1\geq1$. In arity $m$ no
product involving two new generators
occurs, so the corrected arity is
exactly the identity cone and is
contractible.
\end{proof}

Statement~\ref{sa08:C23} also applies to
the low-degree odd parameter introduced
by Statement~\ref{sa08:C20}, provided
the indicated finite-range hypotheses
hold at its new weight. It concerns
one original disk at one fixed arity.
Previously chosen nonlinear repairs of
a whole algebra do not automatically
factor through $D_m(A^{w,q})$.
The proposition therefore does not
assert compatible repairs of mixed
terms, repairs relative to prescribed
boundary data, or one choice of weight
valid for infinitely many arities.
Those distinctions remain necessary
in the integral C.3 problem.
